\section{Trust-Aware Two-Timescale Scheduling}
\label{sec:sched}

This section is the core of \voss{}. It first states the decision problem and its objectives precisely (Sections~\ref{sec:why2}--\ref{sec:objectives}), then says what information each decision uses and how it was obtained (Section~\ref{sec:info}), and then describes the slow planner, the fast placer and the estimators. Table~\ref{tab:notation} lists the notation.

\begin{figure*}[t]
\centering
\resizebox{\textwidth}{!}{% The two loops inside a scheduler. Native width 139 mm.
\begin{tikzpicture}[x=1mm,y=1mm]
  \tikzset{box/.append style={font=\scriptsize}}
  \node[zone, fill=blue!3, minimum width=67mm, minimum height=82mm, anchor=north west] (slow) at (0,0) {};
  \node[zone, fill=green!4, minimum width=67mm, minimum height=82mm, anchor=north west] (fast) at (72,0) {};
  \node[font=\scriptsize\bfseries, anchor=north west] at ([xshift=2mm,yshift=-1mm]slow.north west) {Slow loop: the planner};
  \node[font=\tinyx, anchor=north west] at ([xshift=2mm,yshift=-5mm]slow.north west) {every 30\,s and when things change};
  \node[font=\scriptsize\bfseries, anchor=north west] at ([xshift=2mm,yshift=-1mm]fast.north west) {Fast loop: the placer};
  \node[font=\tinyx, anchor=north west] at ([xshift=2mm,yshift=-5mm]fast.north west) {for every invocation, in milliseconds};
  % slow loop
  \node[box, minimum width=60mm, fill=white] (in) at (33.5,-20) {\textbf{Estimates}\\[-1pt]\tinyx survival, omission rate, assurance,\\[-2pt]\tinyx failure domains, forecast demand};
  \node[box, minimum width=60mm, fill=white] (pl) at (33.5,-37) {\textbf{Robust planner}\\[-1pt]\tinyx greedy or exact; covers demand\\[-2pt]\tinyx even if any $f$ machines fail};
  \node[box, minimum width=60mm, fill=white] (pn) at (33.5,-54) {\textbf{Plan}\\[-1pt]\tinyx who holds which key, which snapshots,\\[-2pt]\tinyx how many warm sandboxes};
  \node[box, minimum width=60mm, fill=white] (ac) at (33.5,-71) {\textbf{Actions to workers}\\[-1pt]\tinyx prepare, pin, warm, release\\[-2pt]\tinyx idempotent, rate-limited};
  % fast loop
  \node[box, minimum width=60mm, fill=white] (iv) at (105.5,-20) {\textbf{Invocation}\\[-1pt]\tinyx tenant, function, deadline, input size};
  \node[box, minimum width=60mm, fill=white] (cd) at (105.5,-37) {\textbf{Candidates}\\[-1pt]\tinyx provisioned machines, plus other\\[-2pt]\tinyx attested ones (needs a key release)};
  \node[box, minimum width=60mm, fill=white] (cm) at (105.5,-54) {\textbf{Cost of each candidate}\\[-1pt]\tinyx wait, trust, start-up, network,\\[-2pt]\tinyx execution, failure risk, exposure};
  \node[box, minimum width=60mm, fill=white] (dc) at (105.5,-71) {\textbf{Decision}\\[-1pt]\tinyx lowest risk-weighted cost;\\[-2pt]\tinyx hedge if risk is high};
  \draw[msg] (in) -- (pl); \draw[msg] (pl) -- (pn); \draw[msg] (pn) -- (ac);
  \draw[msg] (iv) -- (cd); \draw[msg] (cd) -- (cm); \draw[msg] (cm) -- (dc);
  \draw[msg, dashed] (pn.east) -- node[lbl, above, pos=0.5, align=center]{confirmed\\state} (cd.west |- pn.east) -- (cm.west |- pn.east);
  % bottom bar
  \node[att, minimum width=139mm, minimum height=7mm, font=\scriptsize] (wk) at (69.5,-90) {Workers (attested appliances)};
  \draw[msg] (ac.south) -- node[lbl,left,pos=0.45]{actions} (ac.south |- wk.north);
  \draw[msg] (dc.south) -- node[lbl,right,pos=0.45]{dispatch, hedge} (dc.south |- wk.north);
  \draw[msg, dotted] ([xshift=-12mm]wk.north |- wk.north) -- node[lbl,rotate=90,above,pos=0.6]{heartbeats} ([xshift=-12mm]wk.north |- 0,-20) -- (in.east |- 0,-20);
\end{tikzpicture}
}
\caption{The two scheduling timescales within a cluster. The planner decides where provisioned keys and warm state live, using estimates
from the scheduler's view. The placer chooses the worker for each invocation from the planner's confirmed sets, using a per-candidate cost
model. Workers confirm plan actions in heartbeats, so the placer never relies on a host that is not yet warm.}
\label{fig:loops}
\end{figure*}

\subsection{Decision problem and why two timescales}
\label{sec:why2}
\textbf{Setting.} Within one cluster (tens of hosts, Section~\ref{sec:cluster}) the scheduler receives a stream of invocations and must (a)~decide which hosts hold which
attested session state, and (b)~assign every invocation to a host. Its \emph{inputs} are the cluster's state reported by workers, estimates learned from past outcomes, and
configured policy (Section~\ref{sec:info}); its \emph{outputs} are plan actions sent to workers (provision, pin, warm, release) and a dispatch message per invocation.
The two kinds of decision differ in what they control and how fast they must be taken (Table~\ref{tab:levels}).

A single per-invocation scheduler cannot manage attested session state, for three reasons. First, state is \emph{set-valued and long-lived}:
a key released to a host stays for the whole session, so what matters is \emph{how many distinct hosts} hold a tenant's key, which is not a cost that
adds up over assignments. Second, robustness is a \emph{worst case over subsets}: whether demand survives the loss of any $f$ hosts is a property
of the whole provisioned set, not of one placement. Third, establishing state takes seconds to tens of seconds (Fig.~\ref{fig:cost}), far above a
per-request budget, so it must be done ahead of demand. Placement, by contrast, is latency-critical and decomposes into an assignment over a
bounded candidate set. Hence the slow loop optimizes the set-valued, robust, session-scale decisions with an integer program or heuristic, and the
fast loop stays cheap.

\begin{table}[t]
\caption{Notation for a cluster of attested workers $J$.}
\label{tab:notation}
\centering\scriptsize
\setlength{\tabcolsep}{3pt}
\begin{tabular}{@{}p{2.1cm}p{6.1cm}@{}}
\toprule
$t,f,j$ & tenant, function, worker; $t(f)$ is the tenant of $f$ \\
$H,T_{\text{plan}}$ & planning horizon, planning period \\
$y_{tj}\in\{0,1\}$ & worker $j$ is provisioned for tenant $t$ over the horizon \\
$s_{fj}\in\{0,1\}$ & snapshot of $f$ pinned on $j$ (tier L2) \\
$w_{fj}\in\mathbb{Z}_{\ge0}$ & warm sandboxes of $f$ kept on $j$ (tier L3) \\
$c_{tj}\ge 0$ & slots of $j$ earmarked for $t$ (soft cap for the placer) \\
$S_j(H),\,p_j$ & survival probability of $j$'s session over $H$; omission probability \\
$e_j$ & effective availability $S_j(H)(1-p_j)$ \\
$a_j,\,\sigma_t,\,a_{\min,t}$ & assurance of $j$; tenant sensitivity; assurance floor \\
$q_t,\,f_t$ & demand quantile (slots) to cover; number of defecting hosts to tolerate \\
$\mathrm{cap}_j,M_j$ & slots and memory of $j$ \\
$d(j)$ & failure domain of $j$ (ASN and city) \\
\bottomrule
\end{tabular}
\end{table}

\subsection{What is minimised, and why not SLO violations directly}
\label{sec:objectives}
\textbf{The quantity we ultimately care about.} For a synchronous invocation $i$ with deadline $D_i$, an \emph{SLO violation} is the event
$V_i=\{\text{the attempt fails}\}\cup\{\text{end-to-end latency}>D_i\}$. The deadline is $D_i=k_f\,(w_f(\mathit{in}_i)+0.05\,\text{s})$, where $w_f$ is the
reference-core execution time of the function on that input and $k_f=5$ (Section~\ref{sec:protocol}). We \emph{report} violations in every experiment (deadline-met fraction,
success, P99). The scheduler's two loops do not minimise them in the same way (Table~\ref{tab:levels}).

\begin{table*}[t]
\caption{What each decision-maker minimises, subject to what, and whether that is SLO violations. Defaults are in Table~\ref{tab:defaults}.}
\label{tab:levels}
\centering\scriptsize
\setlength{\tabcolsep}{3pt}
\begin{tabular}{@{}p{1.5cm}p{2.0cm}p{4.8cm}p{4.2cm}p{3.2cm}@{}}
\toprule
Decision-maker & Decides & Minimises & Subject to & SLO violations? \\
\midrule
Fast loop (placer) & host for each invocation; hedge; admit or reject & $J_{ij}=w_i P^{\text{miss}}_{ij}+\varepsilon\,\mathbb{E}[T_{ij}]+\lambda_{\exp}x_{ij}$ over hosts $j$ (Eq.~\ref{eq:J}) & eligibility (valid attested session past the expected finish, not draining, assurance $\ge$ floor), top-$K$ candidates, hedge budget & \textbf{Yes, per invocation:} $P^{\text{miss}}$ is that invocation's violation probability; latency and exposure are regularizers \\
Slow loop (planner) & provisioned sets $y$, pinned snapshots $s$, warm VMs $w$, earmarked slots $c$ & uncovered demand, donor memory and disk, key exposure, plan churn, less warm-start benefit (Eq.~\ref{eq:obj}) & robust coverage of the demand quantile against any $f_t$ defectors, slots, memory, disk, failure-domain caps, assurance floor, attestation freshness & \textbf{No, by design:} a surrogate of violation risk (invocations forced onto the unprovisioned path) plus costs SLOs do not see \\
Router & cluster for each request & expected completion time, plus a penalty on clusters whose observed miss rate exceeds $\theta$, plus the price of an unprovisioned start (Section~\ref{sec:cluster}) & cluster not excluded; prefers free slots & Partly: observed miss rate enters as a penalty \\
\bottomrule
\end{tabular}
\end{table*}

\textbf{Fast loop: each invocation's expected violation.} For invocation $i$ of tenant $t$ and candidate host $j$, let $b_i=D_i-(\text{now}-t^{\text{arrive}}_i)$ be the remaining deadline budget and
$T_{ij}$ the predicted completion time (Section~\ref{sec:placer}). The placer minimises
\begin{equation}
\begin{gathered}
J_{ij}=w_i\,P^{\text{miss}}_{ij}+\varepsilon\,\mathbb{E}[T_{ij}]+\lambda_{\exp}\,x_{ij},\\
x_{ij}=\sigma_t(1-a_j)\,\mathbb{1}[j\notin\mathrm{KH}(t)],\\
P^{\text{miss}}_{ij}=p^{\text{fail}}_{ij}+\bigl(1-p^{\text{fail}}_{ij}\bigr)\Pr(T_{ij}>b_i),\\
p^{\text{fail}}_{ij}=1-S_j(\bar e_{ij}+2\,\text{s})(1-p_j).
\end{gathered}
\label{eq:J}
\end{equation}
Here $P^{\text{miss}}_{ij}$ is the probability that invocation $i$ violates its SLO if placed on $j$, with a failure (session end or omission) counted as a violation, and $\bar e_{ij}$ is
the predicted mean execution time. The default weights are $w_i=1$ for synchronous and $0.3$ for asynchronous and filler work, $\varepsilon=0.05\,\text{s}^{-1}$ and $\lambda_{\exp}=1$, so $J$ is measured in
\emph{violation-equivalents}: one second of expected latency is priced like a $5$-point increase in violation probability, and giving a medium-sensitivity tenant's key to a new host of assurance $0.8$ costs
$1\cdot(1-0.8)=0.2$, the same as a 20\% chance of violating the SLO. The placer therefore \emph{does} minimise SLO violations for the invocation in front of it, and latency and key exposure are priced against that
objective, not the other way round. We chose per-invocation violation probability over a P99-based objective because P99 is a population statistic that cannot be attributed to one decision;
minimising each invocation's own miss probability against its own remaining budget is the decision-level quantity whose mean is the deadline-miss rate and whose tail drives the P99. Expectation-, quantile-, CVaR-
and tardiness-based alternatives are compared in an ablation (E03).

\textbf{Slow loop: a surrogate, and why.} The planner decides which hosts hold keys and warm state over the next horizon $H$ ($H=60$\,s; Table~\ref{tab:defaults}). It does \emph{not} minimise violations directly, for five reasons.
\begin{enumerate}
\item \emph{Circularity.} Whether an invocation violates its SLO depends on the placer's later decisions, which depend on the plan; minimising violations would be a bilevel problem solved at every planning round.
\item \emph{No per-request information.} At planning time the invocations do not exist. Only a forecast of demand exists, so the planner can reason about capacity relative to that forecast, not about individual deadlines.
\item \emph{Rare, noisy events.} Violations are 1--5\% events. An objective built from their empirical count has high variance; a coverage requirement at a fixed demand quantile is stable and interpretable.
\item \emph{Violations are not the only cost.} Warming everything everywhere would minimise violations and hold every tenant's key on every host. The scarce resources are donor memory and disk and key exposure
(Section~\ref{sec:exposure}), so the planner minimises cost \emph{under a coverage requirement}, and the trade-off is reported as a Pareto front (Section~\ref{sec:res-pareto}).
\item \emph{No hard guarantee exists.} Under churn and adversaries no plan can guarantee a deadline; coverage against \emph{any} $f_t$ defecting hosts is the strongest statement that assumes no model of the adversary.
\end{enumerate}
\textbf{How the planner still serves the SLO.} The coverage constraint is a chance-constraint surrogate. If invocations of tenant $t$ arrive as a Poisson process of rate $\lambda_t$ with mean service time $s_t$, the number running concurrently is
Poisson with mean $\lambda_t s_t$ (an M/G/$\infty$ system), so $q_t=Q_{0.95}(\lambda_t s_t)$ slots suffice to start every invocation immediately with probability $0.95$ under the forecast. Requiring provisioned,
\emph{available} capacity of at least $q_t$ even after the $f_t$ largest allocations are lost (Eq.~\ref{eq:robust}) bounds the share of invocations forced onto the unprovisioned path, which is the dominant cause of violations on volunteer
hosts. Uncovered demand is priced at $\Lambda_t=10$ per slot, two orders of magnitude above the price of memory ($10^{-4}$ per MB, 0.05 for a 512\,MB VM) and of exposure (about $0.2$ per added key holder), so coverage is bought unless it is infeasible within the
budgets, in which case the shortfall is reported. The Poisson assumption is the planner's \emph{model}: simulated arrivals are a Markov-modulated process with $4\times$ bursts, so realized coverage is below 95\% during bursts. This is a design rationale, not a theorem,
and the realized violation rate is what the experiments measure.

\textbf{What is deliberately not minimised.} There is no explicit P99 target, no throughput objective, no fairness objective across donors and no provider cost term; cost against a commercial reference is reported afterwards (Section~\ref{sec:res-simple}). Reputation never enters
$p^{\text{fail}}$ (Section~\ref{sec:receipts}).

\begin{table}[t]
\caption{Default parameters of the objectives (assumed unless noted; swept in the cited experiments).}
\label{tab:defaults}
\centering\scriptsize
\setlength{\tabcolsep}{3pt}
\begin{tabular}{@{}p{3.1cm}p{5.1cm}@{}}
\toprule
Parameter & Value \\
\midrule
Planning period, horizon & $T_{\text{plan}}=30$\,s, $H=2T_{\text{plan}}=60$\,s \\
Hot set (planned keys) & top 20 tenants, 50 functions, forecast rate $\ge0.01$/s \\
Demand quantile, robustness & $q_t=Q_{0.95}$; $f_t=\lceil0.1\,N_{\text{cluster}}\rceil$ \\
Planner prices & $\Lambda=10$, $\rho_{\text{mem}}=10^{-4}$/MB, $\rho_{\text{ssd}}=2\cdot10^{-5}$/MB, $\lambda=1$, $\kappa_{\text{switch}}=0.05$, $\beta=0.05$, $\gamma=0.5$ \\
Failure-domain cap & $\kappa_{\text{dom}}=0.5$ \\
Placer weights & $w_{\text{sync}}=1$, $w_{\text{other}}=0.3$, $\varepsilon=0.05$, $\lambda_{\exp}=1$, $\sigma\in\{0.1,1,3\}$ \\
Candidates, hedge & $K=12$; $\theta=0.2$, budget 10\% \\
Deadline & $D=5\,(w_f+0.05\,\text{s})$ for sync functions \\
\bottomrule
\end{tabular}
\end{table}

\subsection{Information model: what each decision knows, and how it knows}
\label{sec:info}
A reviewer should be able to tell, for every number a decision uses, whether it was observed, learned from history, forecast, or assumed. We classify every scheduler input into five kinds.
\textbf{Observed (O)}: reported by an attested agent or trusted router by message, stale by at most the heartbeat period of 2\,s plus network delay, with the scheduler's own dispatches since the last report added.
\textbf{Learned from history (H)}: an estimator updated online from \emph{past outcomes only}. \textbf{Forecast (F)}: a quantity computed from H-quantities for the future.
\textbf{Calibrated (C)}: a constant or formula fixed before the run from the calibration bundle, the function catalogue or the geography model; \emph{not} learned online.
\textbf{Policy (P)}: configured by the tenant or operator. Table~\ref{tab:info} lists every quantity.

\begin{table*}[t]
\caption{Every input to a scheduling decision: its kind (O observed, H learned from history, F forecast, C calibrated constant, P policy), its source and update rule, and its value before data exist.}
\label{tab:info}
\centering\scriptsize
\setlength{\tabcolsep}{2.5pt}
\begin{tabular}{@{}p{3.2cm}p{1.5cm}p{0.9cm}p{7.9cm}p{3.7cm}@{}}
\toprule
Quantity & Used by & Kind & Source and update rule & Before enough data \\
\midrule
Slots, memory and disk budgets, speed, link rates, location, tier, session start & planner, placer & O & declared by the worker agent at join; \textbf{speed and link rates are taken as exact} & -- \\
Busy and queued slots, cache tier per function, provisioned tenants, expiry, drain deadline & planner, placer & O & heartbeat every 2\,s; own dispatches since then counted as in flight; hosts silent for 7\,s are left out of plans & host left out \\
Assurance $a_j$ & planner, placer & O & issued by the verifier: substrate base $\times$ firmware-policy factor & -- \\
Execution time on a host & placer & H & Bayesian log-linear regression per function, forgetting factor $0.995$, per-worker offsets shrunk toward 0 (prior count 5), conformal rescaling of the spread over the last 500 residuals & catalogue prior (reference work, size exponent, $\sigma=0.35$) \\
Wait for a slot & placer & O+H & queue length $(\text{busy}+\text{queued}+\text{in flight}-\text{slots}+1)^+$ times the scheduler's EWMA of observed run times (weights $0.98/0.02$), divided by slots & $1$\,s \\
Start-up and trust-establishment time per tier and substrate & placer, planner & C & means of the substrate profile (calibration bundle); \textbf{not re-learned online} & -- \\
Network RTT and transfer time & placer, router & C & great-circle distance at 100\,km per ms $\times1.3$; bytes over the declared link rate; \textbf{no jitter, no measurement} & -- \\
Session survival $S_j(\tau)$ & planner, placer & H$\to$F & per-tier Weibull by censored maximum likelihood on completed sessions and live ages, refit every 60\,s; hour-of-day hazard multiplier after 48 session ends; drain notices only shorten & shape $0.7$, scale 7200\,s until 20 sessions of the tier ended \\
Omission probability $p_j$ & planner, placer, detectors & H & Beta posterior of accepted invocations that return no result, mid-session only; CUSUM ($k=0.1$, $h=5$) resets a host to Beta$(2,4)$ & Beta$(0.5,20)$, mean $0.024$ \\
Arrival rate per tenant and function & planner & H$\to$F & counts in 10\,s bins; EWMAs with $\tau=60$ and $900$\,s; forecast is the maximum of both and of the regional rate pushed by routers & not planned until seen \\
Demand quantile $q_t$ & planner & F & Poisson 95\% quantile of rate $\times$ service time (M/G/$\infty$); service time per key is an EWMA ($0.9/0.1$) of observed run times & catalogue reference work \\
Warm-start benefit $b_f$ & planner & F+C & rate $\times H\times$ (rootfs fetch + boot + init seconds saved) per warm unit & -- \\
Tenant sensitivity, assurance floor, robustness $f_t$ & all & P & tenant configuration & default $f_t=\lceil0.1\,N\rceil$ \\
Objective weights & planner, placer & P & configuration, assumed, swept in E04, E19, E23 & -- \\
Cluster load, observed miss rate, overhead, cold prices & router & O & scheduler summary every 1\,s; cold prices are what the cluster \emph{measured} & priors $3$\,s and $0.5$\,s \\
\bottomrule
\end{tabular}
\end{table*}

\textbf{What is never used.} The true session end time, the true run time and noise of an invocation, a host's true adversarial behaviour, future arrivals, and the true state of other hosts. The code base forbids the policy layer from importing the simulated world, which enforces this
statically. A baseline that does use ground truth exists for analysis only and is never part of \voss{}.

\textbf{How much is history, how much is forecast, how much is assumed.} The \emph{slow loop is mostly forecast}: demand and survival over a 60\,s horizon, extrapolated from rates smoothed over minutes and from session-length statistics accumulated over the whole run.
The \emph{fast loop is mostly observation}: occupancy, cache state and provisioning are at most a few seconds old, and only the execution-time distribution, survival over the next few seconds and the omission probability are statistical estimates learned from past outcomes.
The \emph{speculative} parts are exactly those marked C in Table~\ref{tab:info}: the scheduler assumes start-up costs and network times equal their calibrated nominal values, and that declared host speeds are exact. A deployment would measure these online; this prototype does not,
which is a stated limitation (Section~\ref{sec:threats}). Sensitivity of results to prediction error or staleness (error-injection sweeps) is not part of the paper's experiments.

\subsection{Slow loop: robust provisioning and warm-state planning}
\label{sec:planner}
\textbf{Inputs.} Every $T_{\text{plan}}=30$\,s the scheduler builds a planning instance from its stale view (Table~\ref{tab:info}). The \emph{pool} is the live hosts that are not draining, whose last heartbeat is under 7\,s old and whose session has not expired, capped at 40 hosts (previous holders first, then slots $\times$ assurance).
The \emph{hot set} is the top 20 tenants and top 50 functions whose forecast rate is at least $0.01$\,s$^{-1}$; the long tail stays on the cold path, which keeps the program small. For each hot tenant the planner receives demand $q_t$ (the Poisson 95\% quantile of forecast rate times service time); for each host the survival $S_j(H)$ over the horizon $H=2T_{\text{plan}}=60$\,s, the omission probability $p_j$, the assurance $a_j$, the failure domain, slots, memory, snapshot space and cache inventory; and tenant policy ($\sigma_t$, $a_{\min,t}$, $f_t$).
All are \emph{estimates} from the scheduler's view; no ground truth is used. \textbf{Output.} A plan $\{y,s,w,c\}$ with $y_{tj}$ (host $j$ holds the key of tenant $t$), $s_{fj}$ (snapshot of $f$ pinned on $j$), $w_{fj}$ (warm sandboxes) and $c_{tj}$ (slots earmarked for $t$ on $j$), valid for $H$.

\textbf{Objective.} The planner minimizes uncovered demand, donor cost, exposure and plan churn, less the value of warm starts:
\begin{multline}
\min \sum_t \Lambda_t u_t + \sum_j \Bigl(\rho_{\text{mem}}\sum_f m_f w_{fj} + \rho_{\text{ssd}}\sum_f z_f s_{fj}\Bigr)\\
+ \lambda \sum_t \sigma_t \sum_j (1-a_j)\,y_{tj}\\
+ \nu \sum_{t,j}\bigl(y_{tj}-y^{\text{prev}}_{tj}\bigr)^{+}
- \beta \sum_{f,j} b_f\,(w_{fj}+\gamma s_{fj}),
\label{eq:obj}
\end{multline}
where $u_t\ge0$ is the uncovered demand of tenant $t$ (priced at the cold-path penalty $\Lambda_t$), $m_f$ and $z_f$ are the memory and
snapshot size of $f$, $b_f$ the start-up time a warm sandbox saves, and the remaining coefficients are prices. Each term is logged separately so
that latency, exposure and donor cost can be traded off explicitly (Section~\ref{sec:res-ft}).

\textbf{Robust coverage.} Provisioned capacity must still cover demand when any $f_t$ hosts defect together (crash, omit, or stall):
\begin{equation}
\begin{gathered}
\sum_j e_j\,c_{tj} - \Bigl(f_t\,\theta_t + \sum_j \pi_{tj}\Bigr) + u_t \ge q_t,\\
\pi_{tj}\ge c_{tj}-\theta_t,\quad \pi_{tj},\theta_t\ge 0\quad\forall j.
\end{gathered}
\label{eq:robust}
\end{equation}
The bracket equals the sum of the $f_t$ largest allocations $c_{tj}$: for non-negative $c$, $\sum_{(f_t)} c = \min_{\theta\ge0}\{f_t\theta+\sum_j(c_{j}-\theta)^+\}$ by LP duality,
the standard reformulation of a cardinality-constrained uncertainty set \cite{bertsimas2004price}. Because $e_j\le1$, removing the $f_t$ largest
\emph{slot allocations} removes at least as much as losing any $f_t$ hosts removes from the coverage sum, so Eq.~\eqref{eq:robust} is a sufficient condition and
remains linear. It makes no probabilistic assumption about adversaries: it holds against $f_t$ hosts with a perfect record, which is what an adversary who
builds a clean history before defecting would be. By default $f_t=\lceil\alpha_{\text{assumed}}N_{\text{cluster}}\rceil$ and tenants may opt down.

\textbf{Remaining constraints.} Allocation requires provisioning ($c_{tj}\le \mathrm{cap}_j\,y_{tj}$, $w_{fj}\le W_{\max}\,y_{t(f)j}$, $s_{fj}\le y_{t(f)j}$). Slots, memory and
snapshot storage of each host are respected. A \emph{failure-domain cap} $\sum_{j\in D}c_{tj}\le \kappa_{\text{dom}}\sum_j c_{tj}$ for every domain $D$ prevents a
correlated outage from removing most of a tenant's capacity. The assurance floor forces $y_{tj}=0$ if $a_j<a_{\min,t}$, and a freshness constraint excludes a
host whose attestation would expire within the horizon without time to re-quote.

\textbf{Solving.} One cluster has tens of workers and a bounded set of hot tenants, so the program has about $10^3$--$3\cdot10^3$ integer variables. We solve it with
HiGHS \cite{huangfu2018highs} under a time limit, warm-started from the previous plan, falling back to the incumbent and then to a greedy
heuristic. The \emph{greedy planner} serves hot tenants in decreasing order of demand times cold penalty. For each tenant it repeatedly adds the
host with the highest $S_j(1-p_j)a_j$ per unit of memory and exposure cost, subject to the assurance floor and domain cap, until coverage with the
$f_t$ largest contributions removed is at least $q_t$; it then adds warm sandboxes and snapshots in decreasing value density until budgets bind. We
adopt a pre-set rule: if the exact program improves P99 by less than 5\% over the heuristic in all regimes, the heuristic is the method and the program is
a reference (Section~\ref{sec:res-plan}).

\textbf{When the plan is recomputed.} The plan is re-solved only if the instance changed materially: the hot set or the host pool differs, or $\ln(1+q_t)$ or the per-function warm-unit demand moved by more than a tolerance of $0.2$ in any bucket. Otherwise the previous plan is reused for another horizon, and only unconfirmed actions are re-sent. Solve time and status are logged (\texttt{solver}, \texttt{plans}), and the modelled decision latency of a re-plan is charged to the clock.

\textbf{From plan to actions.} The plan $\{y,s,w,c\}$ is differenced against its predecessor into idempotent, rate-limited actions sent to workers:
\texttt{PREPARE} (obtain the key from the KBS, fetch and decrypt bundles), \texttt{PIN\_SNAPSHOT}, \texttt{WARM} and \texttt{RELEASE}. Workers confirm completion in
heartbeats, and the plan is treated as valid only for state they have \emph{confirmed}, so the placer never counts on a host that is not yet warm.

\subsection{Fast loop: invocation placement}
\label{sec:placer}
\textbf{Inputs and output.} For each invocation $i$ the placer receives only metadata (function, tenant, deadline budget $b_i$, input size, idempotence, SLO class), the scheduler's cluster view (at most one heartbeat old plus its own dispatches) and the current plan. It outputs one of: a host $j$ (optionally with a hedge host), \emph{reject} or \emph{defer}.

\textbf{Candidates and feasibility.} Candidates are the best $K=12$ hosts by a cheap lower bound, drawn from (a)~confirmed provisioned hosts and (b)~other attested hosts (the \emph{cold path}). A host is infeasible if it is draining or its session would expire within 5\,s of the predicted completion, if its assurance is below the tenant floor, or if it lacks memory for a new sandbox.

\textbf{Cost matrix.} For each feasible pair $(i,j)$ the model forms named components, logged per invocation: \emph{wait} for a slot $\bigl(\text{busy}+\text{queued}+\text{in flight}-\text{slots}+1\bigr)^{+}\cdot\bar s/\text{slots}$ with $\bar s$ the EWMA service time; \emph{trust} (zero if $K[t,j]=1$, else key-release round trip plus policy check, calibrated); \emph{start-up} (by the function's tier on $j$ and its disk tier, calibrated); \emph{network} (input and output over the geographic model); and \emph{execution}, a log-normal from the Bayesian regression scaled by declared host speed and substrate slowdown.
The completion time $T_{ij}$ is a moment-matched shifted log-normal over these components. The failure probability is $p^{\text{fail}}_{ij}=1-S_j(\bar e+2\,\text{s})(1-p_j)$, with $\bar e$ the predicted mean execution time: the host must survive the run and not omit. Failure counts as a violation in Eq.~\eqref{eq:J}, so the cost trades a host that is fast but about to leave against one that is slower but safe.

\textbf{Placers.} The default placer assigns each synchronous invocation to the feasible host with the lowest $J_{ij}$ (Eq.~\eqref{eq:J}). For asynchronous and filler work, which arrives in small windows, a minimum-cost flow (\texttt{linear\_sum\_assignment}) maps window invocations to (host, slot) pairs in polynomial time. All baseline placers use the same candidates and cost matrix, so differences are due to the rule alone.

\textbf{Admission, hedging and retries.} \emph{Admission}: if every candidate has $P_{\text{miss}}\ge 0.95$ \emph{and} even the shortest slot wait exceeds the remaining budget (overload), a synchronous request is rejected and an asynchronous one deferred; misses caused only by cold start-up are admitted, because they warm the caches later requests need.
\emph{Hedging}: if the best candidate has $P_{\text{miss}}>\theta=0.2$, the function is idempotent and the cluster's hedge budget (10\% of placements, token bucket) allows, a tied duplicate goes to the best candidate in a \emph{different failure domain}; the loser is cancelled on the first result \cite{dean2013tail}. Attested receipts make duplicates integrity-safe.
\emph{Retries}: on a watchdog or a failure the invocation is re-placed with the failed host excluded, preferring provisioned hosts, within a per-invocation budget and a per-cluster retry token bucket.

\subsection{Estimators}
\label{sec:estimators}
Four estimators feed both loops and are trained online, using only outcomes visible to the scheduler. Parameters are in Table~\ref{tab:info}.
\emph{Survival} $S_j(\tau)$ is a Weibull distribution per tier fitted by censored maximum likelihood on completed sessions and live session ages, refit every 60\,s from a prior (shape $0.7$, scale 7200\,s) until 20 sessions of the tier have ended, with an hour-of-day hazard multiplier once 48 have ended; drain notices only shorten predicted life, since they may be fake.
\emph{Omission} $p_j$ is a Beta posterior over accepted invocations that return a result, counting mid-session outcomes only so that session ends are not double counted; a Beta-CUSUM change-point detector ($k=0.1$, $h=5$) resets a flagged host's eligibility and its posterior to Beta$(2,4)$, which defeats hosts that behave until they defect.
\emph{Execution time} is a Bayesian log-linear regression per function with forgetting factor $0.995$ and shrinkage of per-worker offsets, and a conformal rescaling of the spread over the last 500 residuals \cite{shafer2008conformal}.
\emph{Arrivals} are forecast per hot tenant and function from counts in 10\,s bins by EWMAs with time constants $60$ and $900$\,s; the forecast is the maximum of both and of the regional rate reported by routers.
What these estimators do \emph{not} learn is stated in Section~\ref{sec:info}: start-up and trust costs, network times and declared speeds are calibrated constants. The planner's robustness (Eq.~\eqref{eq:robust}) does not depend on the omission estimator being right; the estimator only reduces how often robustness is exercised.

\subsection{Cluster arrangement and routing}
\label{sec:cluster}
Clusters keep the planner small and intra-cluster latency low. The orchestrator forms them by $k$-medoids over a round-trip-time embedding with $k=\lceil N/N_{\text{target}}\rceil$, requires each to
contain enough failure domains and assurance mass for robust plans to be feasible, and adjusts them on triggers (size bounds, SLO health, demand imbalance, forecast, diversity deficit, scheduler loss) with hysteresis,
dwell times and a migration cost, since a moved worker loses its provisioned state. A router (one per region, trusted for availability) scores each cluster $c$ for an invocation of tenant $t$ and function $f$ by
\begin{equation}
\begin{split}\text{score}(c)={}&\text{rtt}_{c}+o_c+\text{backlog}_c+h_s\max(0,\hat m_c-\theta)\\&+\pi_c(t)\,[\text{no provisioned capacity}]+\text{cold}_c(f),\end{split}
\label{eq:route}
\end{equation}
where $h_s$ is a health weight, $o_c$ the cluster's measured scheduling overhead, $\hat m_c$ its observed miss rate (from its summary every 1\,s), $\pi_c(t)$ the measured key-release price plus an exposure term, and $\text{cold}_c(f)$ the measured cold start. Spillover to another cluster is therefore a priced decision, and a router falls back to the next-best cluster when the dispatch watchdog fires (1\,s).
